\documentclass[11pt]{article}
% Shared preamble for "The Kodamai Container Program: A Complete Technical Record"
% Self-contained; compiles with pdflatex.
\usepackage[utf8]{inputenc}
\usepackage[T1]{fontenc}
\usepackage{lmodern}
\usepackage{amsmath,amssymb,amsthm}
\usepackage{stmaryrd}   % \llbracket \rrbracket
\usepackage{mathrsfs}
\usepackage[margin=1in]{geometry}
\usepackage{enumitem}
\usepackage{microtype}
\usepackage{newunicodechar}
\usepackage[colorlinks=true,linkcolor=blue,citecolor=blue,urlcolor=blue]{hyperref}

% ---- Theorem environments -------------------------------------------------
\theoremstyle{plain}
\newtheorem{theorem}{Theorem}[section]
\newtheorem{proposition}[theorem]{Proposition}
\newtheorem{corollary}[theorem]{Corollary}
\newtheorem{lemma}[theorem]{Lemma}
\theoremstyle{definition}
\newtheorem{definition}[theorem]{Definition}
\newtheorem{result}[theorem]{Result}
\theoremstyle{remark}
\newtheorem*{remark}{Remark}
\newtheorem*{flag}{Flag}

% ---- Grade / registry helper macros ---------------------------------------
% Trust grades are copied verbatim from the master outline; never inflated.
\newcommand{\grade}[1]{\textbf{[#1]}}
\newcommand{\node}[1]{\texttt{#1}}

% ---- Unicode -> LaTeX mapping (so the outline prose transcribes verbatim) --
% Superscripts / subscripts
\newunicodechar{¹}{\ensuremath{^{1}}}
\newunicodechar{²}{\ensuremath{^{2}}}
\newunicodechar{³}{\ensuremath{^{3}}}
\newunicodechar{⁰}{\ensuremath{^{0}}}
\newunicodechar{⁺}{\ensuremath{^{+}}}
\newunicodechar{₁}{\ensuremath{_{1}}}
\newunicodechar{₂}{\ensuremath{_{2}}}
\newunicodechar{₃}{\ensuremath{_{3}}}
\newunicodechar{₄}{\ensuremath{_{4}}}
\newunicodechar{₊}{\ensuremath{_{+}}}
% Greek
\newunicodechar{Γ}{\ensuremath{\Gamma}}
\newunicodechar{Δ}{\ensuremath{\Delta}}
\newunicodechar{Π}{\ensuremath{\Pi}}
\newunicodechar{Σ}{\ensuremath{\Sigma}}
\newunicodechar{Ψ}{\ensuremath{\Psi}}
\newunicodechar{δ}{\ensuremath{\delta}}
\newunicodechar{ε}{\ensuremath{\varepsilon}}
\newunicodechar{θ}{\ensuremath{\theta}}
\newunicodechar{λ}{\ensuremath{\lambda}}
\newunicodechar{ν}{\ensuremath{\nu}}
\newunicodechar{φ}{\ensuremath{\varphi}}
\newunicodechar{ω}{\ensuremath{\omega}}
% Blackboard / script / fraktur
\newunicodechar{ℕ}{\ensuremath{\mathbb{N}}}
\newunicodechar{ℤ}{\ensuremath{\mathbb{Z}}}
\newunicodechar{ℰ}{\ensuremath{\mathcal{E}}}
\newunicodechar{𝒟}{\ensuremath{\mathcal{D}}}
\newunicodechar{𝒯}{\ensuremath{\mathcal{T}}}
\newunicodechar{𝔤}{\ensuremath{\mathfrak{g}}}
\newunicodechar{𝔽}{\ensuremath{\mathbb{F}}}
% Arrows and relations
\newunicodechar{→}{\ensuremath{\to}}
\newunicodechar{↔}{\ensuremath{\leftrightarrow}}
\newunicodechar{⟹}{\ensuremath{\implies}}
\newunicodechar{⟺}{\ensuremath{\iff}}
\newunicodechar{≃}{\ensuremath{\simeq}}
\newunicodechar{≅}{\ensuremath{\cong}}
\newunicodechar{≠}{\ensuremath{\neq}}
\newunicodechar{≡}{\ensuremath{\equiv}}
\newunicodechar{≤}{\ensuremath{\leq}}
\newunicodechar{≥}{\ensuremath{\geq}}
\newunicodechar{⊇}{\ensuremath{\supseteq}}
\newunicodechar{⊊}{\ensuremath{\subsetneq}}
\newunicodechar{∈}{\ensuremath{\in}}
% Operators
\newunicodechar{×}{\ensuremath{\times}}
\newunicodechar{∂}{\ensuremath{\partial}}
\newunicodechar{−}{\ensuremath{-}}
\newunicodechar{∘}{\ensuremath{\circ}}
\newunicodechar{∞}{\ensuremath{\infty}}
\newunicodechar{∧}{\ensuremath{\wedge}}
\newunicodechar{∨}{\ensuremath{\vee}}
\newunicodechar{⊕}{\ensuremath{\oplus}}
\newunicodechar{⊗}{\ensuremath{\otimes}}
\newunicodechar{⊙}{\ensuremath{\odot}}
\newunicodechar{⋆}{\ensuremath{\star}}
\newunicodechar{⋈}{\ensuremath{\bowtie}}
\newunicodechar{⋉}{\ensuremath{\ltimes}}
\newunicodechar{⋊}{\ensuremath{\rtimes}}
\newunicodechar{▷}{\ensuremath{\triangleright}}
\newunicodechar{◁}{\ensuremath{\triangleleft}}
\newunicodechar{⟦}{\ensuremath{\llbracket}}
\newunicodechar{⟧}{\ensuremath{\rrbracket}}
% Punctuation
\newunicodechar{–}{--}
\newunicodechar{—}{---}
\newunicodechar{§}{\S}


\title{The Kodamai Container Program: A Complete Technical Record\\[4pt]
\large Part III --- Differential Structure}
\author{MacBeth}
\date{2026-10-10}

\begin{document}
\maketitle
\begin{center}\itshape Prepared for Neil Ghani.\end{center}

\begin{abstract}
Part~III records the differential structure of the container program. The AAGM one-hole container
derivative $∂$ is the differential combinator of a Cartesian differential category, equivalently a
tangent structure realised by dual-number substitution $T(p)=p(y+εv)\bmod ε^2$. The uniqueness is
peer-reviewed; the Cartan/negation boundary is the honest limit. Every result carries its trust
grade verbatim and its registry node.
\end{abstract}

\paragraph{Grade legend.}
$\grade{peer-reviewed} > \grade{proved} > \grade{lean-verified} / \grade{computed}$. Grades are
copied verbatim and never inflated.

\paragraph{Thesis of Part III.}
The AAGM one-hole container derivative $∂$ is not an ad-hoc operation: it is the differential
combinator of a Cartesian differential category, equivalently a tangent structure realised by
dual-number substitution $T(p)=p(y+εv)\bmod ε^2$. The uniqueness is peer-reviewed; the
Cartan/negation boundary is the honest limit.

\section{The container derivative is a tangent structure}
\emph{Precis.} $\mathrm{SPoly}$ (finite-support polynomial functors, substitution) is a Cartesian
differential category with combinator $D=∂$; hence a Cartesian tangent category.

\begin{theorem}[$∂$ is a CDC differential combinator, \grade{proved}]
$∂$ is the differential combinator of a Cartesian differential category of polynomial functors.
Node \node{tangent-container-derivative}.
\end{theorem}

\begin{theorem}[Dual-number tangent, \grade{proved}]
$T(p)=p(y+εv)\bmod ε^2$ is a strict $\triangleleft$-monoidal base change
$\mathrm{DCont}(k)\to\mathrm{DCont}(D)$, $D=k[ε]/ε^2$, with the laxator equal to the chain rule on
the nose. Node \node{tangent-restricts-to-dcont}.
\end{theorem}

\section{Uniqueness of the container tangent structure}
\emph{Precis.} On finite-support polynomial functors, $∂$ is the unique nontrivial representable
tangent structure. This is the strongest-graded differential result.

\begin{theorem}[Uniqueness, \grade{peer-reviewed}]
On $\mathrm{SPoly}$, $∂$ is the unique nontrivial representable tangent structure, in the
\textbf{corrected statement} per the Rick referee ($2026\text{-}10\text{-}05$). Node
\node{tangent-uniqueness-spoly}.
\end{theorem}
\begin{flag}
State the corrected form, not the earlier unconditional slogan.
\end{flag}

\section{The Cartan / negation boundary}
\emph{Precis.} The full Cartan calculus requires fiberwise negation, which $\mathrm{SPoly}$ lacks;
ring-completion $ℕ\to ℤ$ switches it on. This is the precise edge of the differential story.

\begin{theorem}[Negation obstruction, \grade{proved}]
$\mathrm{SPoly}/∂$ has no fiberwise negation $\implies$ no scalar ring object $\implies$ the Cartan
calculus fails over $ℕ$. Node \node{cartan-scalar-object-spoly}.
\end{theorem}

\begin{theorem}[Ring-completion switches Cartan on, \grade{proved}]
Fiberwise group completion $ℕ\to ℤ$ switches the full Cartan calculus ON ($\mathrm{SPoly}_ℤ$) ---
the far side of the ``negation knife.'' Node \node{virtual-container-cartan}.
\end{theorem}

\begin{theorem}[Lie-parallelizability port, \grade{proved}]
A poly comonoid is tangent-parallelizable $\iff$ an out-degree condition; the Lanfranchi port is
ill-posed / composition-blind. Node \node{lie-parallelizable-poly-comonoid}.
\end{theorem}

\end{document}
